\chapter{Testing and Validation}

\section*{Introduction}

An infrastructure monitoring platform must itself be verifiable and reliable.
This chapter presents the testing strategy adopted for \caimp{}, the results of
unit, integration and load tests, and the traceability matrix demonstrating
functional requirement coverage.

\section{Testing Strategy: The Pyramid}

\caimp{}'s testing strategy follows the classic pyramid, with particular emphasis
on integration tests given the distributed nature of the system.

\begin{figure}[H]
\centering
\begin{tikzpicture}
  \fill[NavyBlue!15, draw=NavyBlue, thick]
    (0,0) -- (6,0) -- (3,4) -- cycle;
  \fill[TealBlue!20, draw=TealBlue, thick]
    (0.8,0) -- (5.2,0) -- (4.0,2.0) -- (2.0,2.0) -- cycle;
  \fill[SuccessGreen!20, draw=SuccessGreen, thick]
    (1.6,0) -- (4.4,0) -- (3.6,1.3) -- (2.4,1.3) -- cycle;
  \node[font=\small\bfseries\color{NavyBlue}]     at (3, 3.2) {E2E / Smoke};
  \node[font=\small\bfseries\color{TealBlue}]     at (3, 1.7) {Integration};
  \node[font=\small\bfseries\color{SuccessGreen}] at (3, 0.65) {Unit};
\end{tikzpicture}
\caption{Testing pyramid for \caimp{}}
\label{fig:pyramid}
\end{figure}

\section{Unit Tests}

Unit tests cover the critical algorithms that can be tested without external
dependencies: anomaly detectors, Holt-Winters engine, change correlation scoring
function, and LLM output validator.

\begin{lstlisting}[language=Python, caption={Unit tests for the Holt-Winters engine}]
class TestHoltWinters:
    def test_stable_series_no_trend(self):
        """Constant series → zero trend, stable forecasts."""
        values = [0.5] * 48
        result = forecast_metric(values, 24, critical_threshold=0.9)
        assert abs(result.trend_slope) < 0.001
        assert all(abs(p["predicted"] - 0.5) < 0.01 for p in result.points)

    def test_rising_series_detects_critical(self):
        """Strongly rising series → time_to_critical is not None."""
        values = [0.5 + 0.01 * i for i in range(48)]  # +1%/h
        result = forecast_metric(values, 24, critical_threshold=0.9)
        assert result.time_to_critical is not None
        assert result.time_to_critical < 24

    def test_confidence_high_for_large_history(self):
        """48+ data points → confidence = 'high'."""
        result = forecast_metric([0.5] * 50, 24, 0.9)
        assert result.confidence == "high"
\end{lstlisting}

\section{Integration Tests}

The smoke test script (\texttt{scripts/smoke\_test\_splunk.py}) implements an
end-to-end test in 7 steps: health checks, authentication, anomaly webhook,
pipeline polling (until \texttt{status = "complete"}), field validation, incident
stats, and runbook CRUD.

\section{Load Tests}

Load test results at 1,000 metrics/second on an 8-core / 16\,GB machine:

\begin{table}[H]
\centering
\caption{Load test results at 1,000 metrics/second}
\label{tab:load-test}
\begin{tabularx}{\textwidth}{lXX}
\toprule
\textbf{Metric} & \textbf{Measured} & \textbf{Target} \\
\midrule
Ingestion latency p50  & 12 ms   & $<$50 ms \\
Ingestion latency p99  & 43 ms   & $<$50 ms \\
Sustained throughput   & 1,180 msg/s & $\geq$1,000 msg/s \\
CPU usage (TW)        & 34\%    & $<$70\% \\
RAM usage (TW)        & 18 MB   & $<$128 MB \\
NATS message loss     & 0       & 0 \\
\bottomrule
\end{tabularx}
\end{table}

All performance objectives defined in Chapter 3 are met.

\section{Security Tests}

\begin{table}[H]
\centering
\caption{Security tests performed}
\label{tab:security-tests}
\small
\begin{tabularx}{\textwidth}{lXl}
\toprule
\textbf{Test} & \textbf{Method} & \textbf{Result} \\
\midrule
SQL injection    & OWASP Top 10 payloads in all text fields & Blocked (parameterised asyncpg) \\
Expired JWT      & Token older than 15 minutes             & HTTP 401 \\
Forged JWT       & Payload modification without re-signing  & HTTP 401 \\
Cross-tenant     & Request with another org's org\_id       & HTTP 404 (RLS) \\
Agent no cert    & OTLP connection without mTLS cert        & TLS handshake failure \\
Revoked cert     & Agent cert in revoked\_agents            & Connection refused \\
\bottomrule
\end{tabularx}
\end{table}

\section{Traceability Matrix}

\begin{table}[H]
\centering
\caption{Functional requirement—test coverage traceability matrix}
\label{tab:traceability-tests}
\small
\begin{tabularx}{\textwidth}{lXll}
\toprule
\textbf{FR} & \textbf{Requirement} & \textbf{Test type} & \textbf{Status} \\
\midrule
FR1  & OTLP collection           & Integration (smoke step 1) & \textcolor{SuccessGreen}{OK} \\
FR6  & 3-algorithm detection     & Unit + Integration          & \textcolor{SuccessGreen}{OK} \\
FR10 & AI explanation            & Integration (smoke step 4)  & \textcolor{SuccessGreen}{OK} \\
FR12 & Change correlation        & Unit (scoring)              & \textcolor{SuccessGreen}{OK} \\
FR14 & Time to critical          & Unit (HW test)              & \textcolor{SuccessGreen}{OK} \\
FR15 & Alert deduplication       & Integration (Redis NX)      & \textcolor{SuccessGreen}{OK} \\
FR17 & Answers-first dashboard   & Integration (endpoint)      & \textcolor{SuccessGreen}{OK} \\
FR18 & Runbook CRUD              & Integration (smoke step 7)  & \textcolor{SuccessGreen}{OK} \\
FR24 & Multi-tenant isolation    & Security (cross-tenant)     & \textcolor{SuccessGreen}{OK} \\
\bottomrule
\end{tabularx}
\end{table}

\section*{Conclusion}

All tests — unit, integration, load, and security — demonstrate the platform's
robustness. Performance targets are met and the principal security vulnerabilities
are covered. The next chapter describes the deployment procedures for putting
the platform into production.
